\ifdefined\gkver\else\def\gkver{student}\fi
\documentclass[\gkver]{gaokaozhenti}
\begin{document}

\chapter{2020年全国理综Ⅰ}
%% paperType: 理综物理部分

\begin{enumerate}
\item
%% number: 1
%% paperName: 2020年全国理综Ⅰ第 1 题 6 分
%% typeId: 1
%% type: 单选题
%% chapter: 运动和力
%% point: 动量定理
%% method: 无
%% score: 6
%% degree: 950
%% duplicateId: 0
%% body:
行驶中的汽车如果发生剧烈碰撞，车内的安全气囊会被弹出并瞬间充满气体。若碰撞后汽车的速度在很短时间内减小为零，关于安全气囊在此过程中的作用，下列说法正确的是\xzanswer{D}
\fourchoices[answer=D]
{增加了司机单位面积的受力大小}
{减少了碰撞前后司机动量的变化量}
{将司机的动能全部转换成汽车的动能}
{延长了司机的受力时间并增大了司机的受力面积}
%% answer: D
%% memo:
\memoanswer{
A．因安全气囊充气后，受力面积增大，故减小了司机单位面积的受力大小，故 A 错误；

B．有无安全气囊司机初动量和末动量均相同，所以动量的改变量也相同，故 B 错误；

C．因有安全气囊的存在，司机和安全气囊接触后会有一部分动能转化为气体的内能，不能全部转化成汽车的动能，故 C 错误；

D．因为安全气囊充气后面积增大，司机的受力面积也增大，在司机挤压气囊作用过程中由于气囊的缓冲故增加了作用时间，故 D 正确。

故选 D。
}

\item
%% number: 2
%% paperName: 2020年全国理综Ⅰ第 2 题 6 分
%% typeId: 1
%% type: 单选题
%% chapter: 曲线运动
%% point: 万有引力定律
%% method: 无
%% score: 6
%% degree: 850
%% duplicateId: 0
%% body:
火星的质量约为地球质量的 $\frac{1}{10}$，半径约为地球半径的 $\frac{1}{2}$，则同一物体在火星表面与在地球表面受到的引力的比值约为\xzanswer{B}
\fourchoices[answer=B]
{0.2}
{0.4}
{2.0}
{2.5}
%% answer: B
%% memo:
\memoanswer{
设物体质量为 $m$，则在火星表面有

$F_{1}=G\frac{M_{1}m}{R_{1}^{2}}$

在地球表面有

$F_{2}=G\frac{M_{2}m}{R_{2}^{2}}$

由题意知有

$\frac{M_{1}}{M_{2}}=\frac{1}{10}$，$\frac{R_{1}}{R_{2}}=\frac{1}{2}$

故联立以上公式可得

$\frac{F_{1}}{F_{2}}=\frac{M_{1}R_{2}^{2}}{M_{2}R_{1}^{2}}=\frac{1}{10}\times\frac{4}{1}=0.4$

故选 B。
}

\item
%% number: 3
%% paperName: 2020年全国理综Ⅰ第 3 题 6 分
%% typeId: 1
%% type: 单选题
%% chapter: 曲线运动
%% point: 竖直平面圆周运动
%% method: 无
%% score: 6
%% degree: 850
%% duplicateId: 0
%% body:
如图，一同学表演荡秋千。已知秋千的两根绳长均为 $10\Um$，该同学和秋千踏板的总质量约为 $50\Ukg$。绳的质量忽略不计，当该同学荡到秋千支架的正下方时，速度大小为 $8\Ums$，此时每根绳子平均承受的拉力约为\xzanswer{B}
\onepicture[align=right, w=2.3cm]{figs/03.png}
\fourchoices[answer=B]
{200 N}
{400 N}
{600 N}
{800 N}
%% answer: B
%% memo:
\memoanswer{
在最低点由 $2T-mg=\frac{mv^{2}}{r}$，可知 $T=410\UN$，即每根绳子拉力约为 $410\UN$，故选 B。
}

\item
%% number: 4
%% paperName: 2020年全国理综Ⅰ第 4 题 6 分
%% typeId: 1
%% type: 单选题
%% chapter: 电路
%% point: 电容
%% method: 无
%% score: 6
%% degree: 750
%% duplicateId: 0
%% body:
图（a）所示的电路中，K 与 L 间接一智能电源，用以控制电容器 C 两端的电压 $U_{C}$。如果 $U_{C}$ 随时间 $t$ 的变化如图（b）所示，则下列描述电阻 $R$ 两端电压 $U_{R}$ 随时间 $t$ 变化的图像中，正确的是\xzanswer{A}
\onepicture[w=11.0cm]{\includesvg{figs/04}}
\fourchoices[answer=A, ispicture=true, h=2.0cm]
{\includesvg{figs/04a}}
{\includesvg{figs/04b}}
{\includesvg{figs/04c}}
{\includesvg{figs/04d}}
%% answer: A
%% memo:
\memoanswer{
根据电容器的定义式 $C=\frac{Q}{U}$ 可知

$U_{C}=\frac{Q}{C}=\frac{I}{C}t$

结合图像可知，图像的斜率为 $\frac{I}{C}$，则 $1\sim2\Us$ 内的电流 $I_{12}$ 与 $3\sim5\Us$ 内的电流 $I_{35}$ 关系为

$I_{12}=2I_{35}$

且两段时间中的电流方向相反，根据欧姆定律 $I=\frac{U}{R}$ 可知 $R$ 两端电压大小关系满足

$U_{R12}=2U_{R35}$

由于电流方向不同，所以电压方向不同。

故选 A。
}

\item
%% number: 5
%% paperName: 2020年全国理综Ⅰ第 5 题 6 分
%% typeId: 1
%% type: 单选题
%% chapter: 磁场
%% point: 带电粒子在磁场中的运动
%% method: 无
%% score: 6
%% degree: 450
%% duplicateId: 0
%% body:
一匀强磁场的磁感应强度大小为 $B$，方向垂直于纸面向外，其边界如图中虚线所示，ab 为半圆，ac、bd 与直径 ab 共线，ac 间的距离等于半圆的半径。一束质量为 $m$、电荷量为 $q$（$q>0$）的粒子，在纸面内从 c 点垂直于 ac 射入磁场，这些粒子具有各种速率。不计粒子之间的相互作用。在磁场中运动时间最长的粒子，其运动时间为\xzanswer{C}
\onepicture[w=7.0cm]{\includesvg{figs/05}}
\fourchoices[answer=C]
{$\frac{7\pi m}{6qB}$}
{$\frac{5\pi m}{4qB}$}
{$\frac{4\pi m}{3qB}$}
{$\frac{3\pi m}{2qB}$}
%% answer: C
%% memo:
\memoanswer{
粒子在磁场中做匀速圆周运动的周期 $T=\frac{2\pi m}{qB}$，粒子在磁场中运动的时间 $t=\frac{\theta}{2\pi}\cdot T=\frac{\theta m}{qB}$

则粒子在磁场中运动的时间与速度无关，轨迹对应的圆心角越大，运动时间越长。采用放缩圆解决该问题，粒子垂直 $ac$ 射入磁场，则轨迹圆心必在 $ac$ 直线上，将粒子的轨迹半径由零逐渐放大。

当半径 $r\leq0.5R$ 和 $r\geq1.5R$ 时，粒子分别从 $ac$、$bd$ 区域射出，磁场中的轨迹为半圆，运动时间等于半个周期。

当 $0.5R<r<1.5R$ 时，粒子从半圆边界射出，逐渐将轨迹半径从 $0.5R$ 逐渐放大，粒子射出位置从半圆顶端向下移动，轨迹圆心角从 $\pi$ 逐渐增大，当轨迹半径为 $R$ 时，轨迹圆心角最大，然后再增大轨迹半径，轨迹圆心角减小，因此当轨迹半径等于 $R$ 时轨迹圆心角最大，即轨迹对应的最大圆心角 $\theta=\pi+\frac{\pi}{3}=\frac{4}{3}\pi$

粒子运动最长时间为 $t=\frac{\theta}{2\pi}T=\frac{\frac{4}{3}\pi}{2\pi}\times\frac{2\pi m}{qB}=\frac{4\pi m}{3qB}$，

故选 C。

\twopicture[showlabel=false, hA=3.6cm, hB=3.6cm]{figs/05_an1.jpg}{figs/05_an.jpg}
}

\item
%% number: 6
%% paperName: 2020年全国理综Ⅰ第 6 题 6 分
%% typeId: 2
%% type: 多选题
%% chapter: 原子和原子核
%% point: 人工核反应
%% method: 无
%% score: 6
%% degree: 850
%% duplicateId: 0
%% body:
下列核反应方程中，$X_{1}$，$X_{2}$，$X_{3}$，$X_{4}$ 代表 $\alpha$ 粒子的有\xzanswer{BD}
\fourchoices[answer=BD]
{\ce{^{2}_{1}H + ^{2}_{1}H -> ^{1}_{0}n} + $X_{1}$}
{\ce{^{2}_{1}H + ^{3}_{1}H -> ^{1}_{0}n} + $X_{2}$}
{\ce{^{235}_{92}U + ^{1}_{0}n -> ^{144}_{56}Ba + ^{89}_{36}Kr} + 3$X_{3}$}
{\ce{^{1}_{0}n + ^{6}_{3}Li -> ^{3}_{1}H} + $X_{4}$}
%% answer: BD
%% memo:
\memoanswer{
$\alpha$ 粒子为氦原子核 $\ce{^{4}_{2}He}$，根据核反应方程遵守电荷数守恒和质量数守恒，A 选项中的 $X_{1}$ 为 $\ce{^{3}_{2}He}$，B 选项中的 $X_{2}$ 为 $\ce{^{4}_{2}He}$，C 选项中的 $X_{3}$ 为中子 $\ce{^{1}_{0}n}$，D 选项中的 $X_{4}$ 为 $\ce{^{4}_{2}He}$。故选 BD。
}

\item
%% number: 7
%% paperName: 2020年全国理综Ⅰ第 7 题 6 分
%% typeId: 2
%% type: 多选题
%% chapter: 功和能
%% point: 功能中的图象
%% method: 无
%% score: 6
%% degree: 550
%% duplicateId: 0
%% body:
一物块在高 $3.0\Um$、长 $5.0\Um$ 的斜面顶端从静止开始沿斜面下滑，其重力势能和动能随下滑距离 $s$ 的变化如图中直线Ⅰ、Ⅱ 所示，重力加速度取 $10\Umsq$。则\xzanswer{AB}
\onepicture[w=6.5cm]{\includesvg{figs/07}}
\fourchoices[answer=AB]
{物块下滑过程中机械能不守恒}
{物块与斜面间的动摩擦因数为 0.5}
{物块下滑时加速度的大小为 $6.0\Umsq$}
{当物块下滑 $2.0\Um$ 时机械能损失了 $12\UJ$}
%% answer: AB
%% memo:
\memoanswer{
A．下滑 $5\Um$ 的过程中，重力势能减少 $30\UJ$，动能增加 $10\UJ$，减小的重力势能并不等于增加的动能，所以机械能不守恒，A 正确；

B．斜面高 $3\Um$、长 $5\Um$，则斜面倾角为 $\theta=37^\circ$。令斜面底端为零势面，则物块在斜面顶端时的重力势能 $mgh=30\UJ$，可得质量 $m=1\Ukg$。

下滑 $5\Um$ 过程中，由功能原理，机械能的减少量等于克服摩擦力做的功

$\mu mg\cos\theta\cdot s=20\UJ$

求得 $\mu=0.5$。B 正确；

C．由牛顿第二定律 $mg\sin\theta-\mu mg\cos\theta=ma$，求得 $a=2\Umsq$。C 错误；

D．物块下滑 $2.0\Um$ 时，重力势能减少 $12\UJ$，动能增加 $4\UJ$，所以机械能损失了 $8\UJ$，D 选项错误。

故选 AB。
}

\item
%% number: 8
%% paperName: 2020年全国理综Ⅰ第 8 题 6 分
%% typeId: 2
%% type: 多选题
%% chapter: 电磁感应
%% point: 双杆问题
%% method: 无
%% score: 6
%% degree: 350
%% duplicateId: 0
%% body:
如图，U 形光滑金属框 abcd 置于水平绝缘平台上，ab 和 dc 边平行，和 bc 边垂直。ab、dc 足够长，整个金属框电阻可忽略。一根具有一定电阻的导体棒 MN 置于金属框上，用水平恒力 $F$ 向右拉动金属框，运动过程中，装置始终处于竖直向下的匀强磁场中，MN 与金属框保持良好接触，且与 bc 边保持平行。经过一段时间后\xzanswer{BC}
\onepicture[w=7.0cm]{figs/08.png}
\fourchoices[answer=BC]
{金属框的速度大小趋于恒定值}
{金属框的加速度大小趋于恒定值}
{导体棒所受安培力的大小趋于恒定值}
{导体棒到金属框 bc 边的距离趋于恒定值}
%% answer: BC
%% memo:
\memoanswer{
由 $bc$ 边切割磁感线产生电动势，形成电流，使得导体棒 $MN$ 受到向右的安培力，做加速运动，$bc$ 边受到向左的安培力，向右做加速运动。当 $MN$ 运动时，金属框的 $bc$ 边和导体棒 $MN$ 一起切割磁感线，设导体棒 $MN$ 和金属框的速度分别为 $v_{1}$、$v_{2}$，则电路中的电动势 $E=BL(v_{2}-v_{1})$

电流中的电流 $I=\frac{E}{R}=\frac{BL(v_{2}-v_{1})}{R}$

金属框和导体棒 $MN$ 受到的安培力 $F_{\text{安框}}=\frac{B^{2}L^{2}(v_{2}-v_{1})}{R}$，与运动方向相反；$F_{\text{安}MN}=\frac{B^{2}L^{2}(v_{2}-v_{1})}{R}$，与运动方向相同

设导体棒 $MN$ 和金属框的质量分别为 $m_{1}$、$m_{2}$，则对导体棒 $MN$

$\frac{B^{2}L^{2}(v_{2}-v_{1})}{R}=m_{1}a_{1}$

对金属框

$F-\frac{B^{2}L^{2}(v_{2}-v_{1})}{R}=m_{2}a_{2}$

初始速度均为零，则 $a_{1}$ 从零开始逐渐增加，$a_{2}$ 从 $\frac{F}{m_{2}}$ 开始逐渐减小。当 $a_{1}=a_{2}$ 时，相对速度

$v_{2}-v_{1}=\frac{FRm_{1}}{2B^{2}L^{2}(m_{1}+m_{2})}$

大小恒定。整个运动过程用速度时间图象描述如下。

\onepicture[w=7.0cm]{figs/08_an.jpg}

综上可得，金属框的加速度趋于恒定值，安培力也趋于恒定值，BC 选项正确；

金属框的速度会一直增大，导体棒到金属框 $bc$ 边的距离也会一直增大，AD 选项错误。

故选 BC。
}

\item
%% number: 9
%% paperName: 2020年全国理综Ⅰ第 9 题 6 分
%% typeId: 4
%% type: 实验题
%% chapter: 电路
%% point: 伏安法测电阻
%% method: 无
%% score: 6
%% degree: 550
%% duplicateId: 0
%% body:
某同学用伏安法测量一阻值为几十欧姆的电阻 $R_{x}$，所用电压表的内阻为 $1\UkO$，电流表内阻为 $0.5\UO$。该同学采用两种测量方案，一种是将电压表跨接在图（a）所示电路的 O、P 两点之间，另一种是跨接在 O、Q 两点之间。测量得到如图（b）所示的两条 $U-I$ 图线，其中 $U$ 与 $I$ 分别为电压表和电流表的示数。

\onepicture[align=right, w=3.6cm]{figs/09a.png}

\onepicture[num=图（b）, labelprefix={}, w=7.4cm]{figs/09b.png}

回答下列问题：
\begin{enumerate}
	\item
	图（b）中标记为 II 的图线是采用电压表跨接在\tkanswer[1.8cm]{O、P}（填“O、P”或“O、Q”）两点的方案测量得到的。

	\item
	根据所用实验器材和图（b）可判断，由图线\tkanswer[1.2cm]{I}（填“I”或“II”）得到的结果更接近待测电阻的真实值，结果为\tkanswer[1.6cm]{50.5}$\UO$（保留 1 位小数）。

	\item
	考虑到实验中电表内阻的影响，需对（2）中得到的结果进行修正，修正后待测电阻的阻值为\tkanswer[1.6cm]{50.0}$\UO$（保留 1 位小数）。
\end{enumerate}
%% answer:
\jdanswer{
\begin{enumerate}
	\item
	O、P

	\item
	I，$50.5\UO$

	\item
	$50.0\UO$

\end{enumerate}
}
%% memo:
\memoanswer{
（1）若将电压表接在 O、P 之间，

则 $U=\frac{R_{x}R_{V}}{R_{x}+R_{V}}\cdot I$

根据一次函数关系可知对应斜率为 $\frac{R_{x}R_{V}}{R_{x}+R_{V}}$。

若将电压表接在 O、Q 之间，电流表分压为 $U_{A}=IR_{A}$，根据欧姆定律变形可知 $R=\frac{U-IR_{A}}{I}$，解得：

$U=I(R+R_{\mathrm{A}})$

根据一次函数可知对应斜率为 $R+R_{A}$，对比图像的斜率可知 $k_{1}>k_{II}$

所以 II 图线是采用电压表跨接在 O、P 之间。

（2）因为待测电阻为几十欧姆的电阻，通过图像斜率大致估算待测电阻为 $50\UO$ 左右，根据

$\frac{1\UkO}{50\UO}<\frac{50\UO}{0.5\UO}$

说明电流表的分压较小，电流表的分流较大，所以电压表应跨接在 O、Q 之间，所以选择图线 I 得到的结果较为准确。

根据图像可知

$R_{x}=\frac{3\UV-1\UV}{59.6\UmA-20\UmA}\approx50.5\UO$

（3）考虑电流表内阻，则修正后的电阻为

$R_{\mathrm{x}}^{\prime}=R_{x}-r_{\mathrm{A}}=50.5\UO-0.5\UO=50.0\UO$
}

\item
%% number: 10
%% paperName: 2020年全国理综Ⅰ第 10 题 9 分
%% typeId: 4
%% type: 实验题
%% chapter: 运动和力
%% point: 动量守恒定律
%% method: 无
%% score: 9
%% degree: 450
%% duplicateId: 0
%% body:
某同学用如图所示的实验装置验证动量定理，所用器材包括：气垫导轨、滑块（上方安装有宽度为 $d$ 的遮光片）、两个与计算机相连接的光电门、砝码盘和砝码等。

\onepicture[align=right, w=8.0cm]{figs/10.png}

实验步骤如下：
\begin{steps}
	\item
	开动气泵，调节气垫导轨，轻推滑块，当滑块上的遮光片经过两个光电门的遮光时间\tkanswer[1.6cm]{大约相等}时，可认为气垫导轨水平；
	\item
	用天平测砝码与砝码盘的总质量 $m_{1}$、滑块（含遮光片）的质量 $m_{2}$；
	\item
	用细线跨过轻质定滑轮将滑块与砝码盘连接，并让细线水平拉动滑块；
	\item
	令滑块在砝码和砝码盘的拉动下从左边开始运动，和计算机连接的光电门能测量出遮光片经过 A、B 两处的光电门的遮光时间 $\Delta t_{1}$、$\Delta t_{2}$ 及遮光片从 A 运动到 B 所用的时间 $t_{12}$；
	\item
	在遮光片随滑块从 A 运动到 B 的过程中，如果将砝码和砝码盘所受重力视为滑块所受拉力，拉力冲量的大小 $I=$\tkanswer[1.8cm]{$m_{1}gt_{12}$}，滑块动量改变量的大小 $\Delta p=$\tkanswer[2.8cm]{$m_{2}\left(\frac{d}{\Delta t_{2}}-\frac{d}{\Delta t_{1}}\right)$}；（用题中给出的物理量及重力加速度 $g$ 表示）
	\item
	某次测量得到的一组数据为：$d=1.000\Ucm$，$m_{1}=1.50\times10^{-2}\Ukg$，$m_{2}=0.400\Ukg$，$\Delta t_{1}=3.900\times10^{-2}\Us$，$\Delta t_{2}=1.270\times10^{-2}\Us$，$t_{12}=1.50\Us$，取 $g=9.80\Umsq$。计算可得 $I=$\tkanswer[1.5cm]{0.221}$\UNs$，$\Delta p=$\tkanswer[1.5cm]{0.212}$\Ukgms$；（结果均保留 3 位有效数字）
	\item
	定义 $\delta=\left|\frac{I-\Delta p}{I}\right|\times100\%$，本次实验 $\delta=$\tkanswer[1.0cm]{4}\%（保留 1 位有效数字）。
\end{steps}
%% answer:
%% memo:
\memoanswer{
（1）当经过 A，B 两个光电门时间相等时，速度相等，此时由于阻力很小，可以认为导轨是水平的。

（5）由 $I=Ft$，知 $I=m_{1}gt_{12}$

由 $\Delta p=mv_{2}-mv_{1}$ 知

$\Delta p=m_{2}\cdot\frac{d}{\Delta t_{2}}-m_{2}\cdot\frac{d}{\Delta t_{1}}=m_{2}\left(\frac{d}{\Delta t_{2}}-\frac{d}{\Delta t_{1}}\right)$

（6）代入数值知，冲量 $I=m_{1}gt_{12}=1.5\times10^{-2}\times9.8\times1.5\UNs=0.221\UNs$

$\Delta p=m_{2}\left(\frac{d}{\Delta t_{2}}-\frac{d}{\Delta t_{1}}\right)=0.212\Ukgms$

动量改变量

（7）$\delta=\frac{|I-\Delta p|}{I}\times100\%=\frac{0.225-0.212}{0.225}\times100\%\approx4\%$
}

\item
%% number: 11
%% paperName: 2020年全国理综Ⅰ第 11 题 12 分
%% typeId: 6
%% type: 计算题
%% chapter: 运动和力
%% point: 牛顿定律的应用
%% method: 无
%% score: 12
%% degree: 650
%% duplicateId: 0
%% body:
我国自主研制了运-20 重型运输机。飞机获得的升力大小 $F$ 可用 $F=kv^{2}$ 描写，$k$ 为系数；$v$ 是飞机在平直跑道上的滑行速度，$F$ 与飞机所受重力相等时的 $v$ 称为飞机的起飞离地速度，已知飞机质量为 $1.21\times10^{5}\Ukg$ 时，起飞离地速度为 $66\Ums$；装载货物后质量为 $1.69\times10^{5}\Ukg$，装载货物前后起飞离地时的 $k$ 值可视为不变。
\begin{enumerate}
	\item
	求飞机装载货物后的起飞离地速度；
	\item
	若该飞机装载货物后，从静止开始匀加速滑行 $1521\Um$ 起飞离地，求飞机在滑行过程中加速度的大小和所用的时间。
\end{enumerate}
%% answer:
\jdanswer{
\begin{enumerate}
	\item
	$v_{2}=78\Ums$

	\item
	$2\Umsq$，$t=39\Us$

\end{enumerate}
}
%% memo:
\memoanswer{
（1）空载起飞时，升力正好等于重力：$kv_{1}^{2}=m_{1}g$

满载起飞时，升力正好等于重力：$kv_{2}^{2}=m_{2}g$

由上两式解得：$v_{2}=78\Ums$

（2）满载货物的飞机做初速度为零的匀加速直线运动，所以 $v_{2}^{2}-0=2ax$

解得：$a=2\Umsq$

由加速的定义式变形得：$t=\frac{v_{2}}{a}$

解得：$t=39\Us$
}

\item
%% number: 12
%% paperName: 2020年全国理综Ⅰ第 12 题 20 分
%% typeId: 6
%% type: 计算题
%% chapter: 磁场
%% point: 带电粒子在磁场中的运动
%% method: 无
%% score: 20
%% degree: 450
%% duplicateId: 0
%% body:
在一柱形区域内有匀强电场，柱的横截面积是以 O 为圆心，半径为 $R$ 的圆，AB 为圆的直径，如图所示。质量为 $m$，电荷量为 $q$（$q>0$）的带电粒子在纸面内自 A 点先后以不同的速度进入电场，速度方向与电场的方向垂直。已知刚进入电场时速度为零的粒子，自圆周上的 C 点以速率 $v_{0}$ 穿出电场，AC 与 AB 的夹角 $\theta=60^\circ$。运动中粒子仅受电场力作用。
\begin{enumerate}
	\item
	求电场强度的大小；
	\item
	为使粒子穿过电场后的动能增量最大，该粒子进入电场时的速度应为多大？
	\item
	为使粒子穿过电场前后动量变化量的大小为 $mv_{0}$，该粒子进入电场时的速度应为多大？
\end{enumerate}
\onepicture[align=right, w=5.5cm]{\includesvg{figs/12}}
%% answer:
\jdanswer{
\begin{enumerate}
	\item
	$E=\frac{mv_{0}^{2}}{2qR}$

	\item
	$v_{1}=\frac{\sqrt{2}v_{0}}{4}$

	\item
	$0$ 或 $v_{2}=\frac{\sqrt{3}v_{0}}{2}$

\end{enumerate}
}
%% memo:
\memoanswer{
（1）由题意知在 A 点速度为零的粒子会沿着电场线方向运动，由于 $q>0$，故电场线由 A 指向 C，根据几何关系可知：

$x_{AC}=R$

所以根据动能定理有：

$qEx_{AC}=\frac{1}{2}mv_{0}^{2}-0$

解得：$E=\frac{mv_{0}^{2}}{2qR}$

（2）根据题意可知要使粒子动能增量最大则沿电场线方向移动距离最多，做 AC 垂线并且与圆相切，切点为 D，即粒子要从 D 点射出时沿电场线方向移动距离最多，粒子在电场中做类平抛运动，根据几何关系有

$x=R\sin60^\circ=v_{1}t$

$y=R+R\cos60^\circ=\frac{1}{2}at^{2}$

而电场力提供加速度有 $qE=ma$

联立各式解得粒子进入电场时的速度：$v_{1}=\frac{\sqrt{2}v_{0}}{4}$

\onepicture[w=6.0cm]{figs/12_an.jpg}

（3）因为粒子在电场中做类平抛运动，粒子穿过电场前后动量变化量大小为 $mv_{0}$，即在电场方向上速度变化为 $v_{0}$，过 C 点做 AC 垂线会与圆周交于 B 点，故由题意可知粒子会从 C 点或 B 点射出。当从 B 点射出时由几何关系有

$x_{BC}=\sqrt{3}R=v_{2}t_{2}$

$x_{AC}=R=\frac{1}{2}at_{2}^{2}$

电场力提供加速度有 $qE=ma$

联立解得 $v_{2}=\frac{\sqrt{3}v_{0}}{2}$；当粒子从 C 点射出时初速度为 0。

另解：由题意知，初速度为 0 时，动量增量的大小为 $mv_{0}$，此即问题的一个解。自 A 点以不同的速率垂直于电场方向射入电场的粒子，动量变化都相同，自 B 点射出电场的粒子，其动量变化量也恒为 $mv_{0}$，由几何关系及运动学规律可得，此时入射速率为 $v_{2}=\frac{\sqrt{3}v_{0}}{2}$。
}

\item
%% number: 13
%% paperName: 2020年全国理综Ⅰ第 13 题 10 分
%% typeId: 6
%% type: 计算题
%% chapter: 气体性质
%% point: 玻意耳定律
%% method: 无
%% score: 10
%% degree: 550
%% duplicateId: 0
%% body:
\textbf{（1）}（5 分）分子间作用力 $F$ 与分子间距 $r$ 的关系如图所示，$r=r_{1}$ 时，$F=0$。分子间势能由 $r$ 决定，规定两分子相距无穷远时分子间的势能为零。若一分子固定于原点 O，另一分子从距 O 点很远处向 O 点运动，在两分子间距减小到 $r_{2}$ 的过程中，势能\tkanswer[1.2cm]{减小}（填“减小”“不变”或“增大”）；在间距由 $r_{2}$ 减小到 $r_{1}$ 的过程中，势能\tkanswer[1.2cm]{减小}（填“减小”“不变”或“增大”）；在间距等于 $r_{1}$ 处，势能\tkanswer[1.2cm]{小于}（填“大于”“等于”或“小于”）零。

\onepicture[align=right, w=5.5cm]{\includesvg{figs/13}}

\textbf{（2）}（10 分）甲、乙两个储气罐储存有同种气体（可视为理想气体）。甲罐的容积为 $V$，罐中气体的压强为 $p$；乙罐的容积为 $2V$，罐中气体的压强为 $\frac{1}{2}p$。现通过连接两罐的细管把甲罐中的部分气体调配到乙罐中去，两罐中气体温度相同且在调配过程中保持不变，调配后两罐中气体的压强相等。求调配后：
\begin{enumerate}
	\item
	两罐中气体的压强；
	\item
	甲罐中气体的质量与甲罐中原有气体的质量之比。
\end{enumerate}
%% answer:
\jdanswer{
\begin{enumerate}
	\item
	$\frac{2}{3}p$

	\item
	$\frac{2}{3}$

\end{enumerate}
}
%% memo:
\memoanswer{
（1）从距 O 点很远处向 O 点运动，两分子间距减小到 $r_{2}$ 的过程中，分子间体现引力，引力做正功，分子势能减小；

\onepicture[w=5.5cm]{\includesvg{figs/13_an}}

在 $r_{2}\to r_{1}$ 的过程中，分子间仍然体现引力，引力做正功，分子势能减小；

在间距等于 $r_{1}$ 之前，分子势能一直减小，取无穷远处分子间势能为零，则在 $r_{1}$ 处分子势能小于零。

（2）(i) 气体发生等温变化，对甲乙中的气体，可认为甲中原气体有体积 $V$ 变成 $3V$，乙中原气体体积有 $2V$ 变成 $3V$，则根据玻意耳定律分别有

$pV=p_{1}\cdot3V$，$\frac{1}{2}p\cdot2V=p_{2}\cdot3V$

则

$pV+\frac{1}{2}p\cdot2V=(p_{1}+p_{2})\times3V$

则甲乙中气体最终压强

$p'=p_{1}+p_{2}=\frac{2}{3}p$

(ii) 若调配后将甲气体再等温压缩到气体原来的压强为 $p$，则 $p'V=pV'$

计算可得

$V'=\frac{2}{3}V$

由密度定律可得，质量之比等于

$\frac{m_{\text{现}}}{m_{\text{原}}}=\frac{V'}{V}=\frac{2}{3}$
}

\item
%% number: 14
%% paperName: 2020年全国理综Ⅰ第 14 题 10 分
%% typeId: 6
%% type: 计算题
%% chapter: 机械波
%% point: 干涉衍射
%% method: 无
%% score: 10
%% degree: 550
%% duplicateId: 0
%% body:
\textbf{（1）}（5 分）在下列现象中，可以用多普勒效应解释的有\xzanswer{BCE}
\fivechoices[answer=BCE]
{雷雨天看到闪电后，稍过一会儿才能听到雷声}
{超声波被血管中的血流反射后，探测器接收到的超声波频率发生变化}
{观察者听到远去的列车发出的汽笛声，音调会变低}
{同一声源发出的声波，在空气和水中传播的速度不同}
{天文学上观察到双星（相距较近、均绕它们连线上某点做圆周运动的两颗恒星）光谱随时间的周期性变化}

\textbf{（2）}（10 分）一振动片以频率 $f$ 做简谐振动时，固定在振动片上的两根细杆同步周期性地触动水面上 a、b 两点，两波源发出的波在水面上形成稳定的干涉图样。c 是水面上的一点，a、b、c 间的距离均为 $l$，如图所示。已知除 c 点外，在 ac 连线上还有其他振幅极大的点，其中距 c 最近的点到 c 的距离为 $\frac{3}{8}l$。求：
\begin{enumerate}
	\item
	波的波长；
	\item
	波的传播速度。
\end{enumerate}
\onepicture[align=right, w=5.5cm]{\includesvg{figs/14}}
%% answer:
\jdanswer{
\begin{enumerate}
	\item
	$\frac{1}{4}l$

	\item
	$\frac{1}{4}fl$

\end{enumerate}
}
%% memo:
\memoanswer{
（1）A．之所以不能同时观察到是因为声音的传播速度比光的传播速度慢，所以 A 错误；

B．超声波与血液中的血小板等细胞发生反射时，由于血小板的运动会使得反射声波的频率发生变化，B 正确；

C．列车和人的位置相对变化了，所以听得的声音频率发生了变化，所以 C 正确；

D．波动传播速度不一样是由于波的频率不一样导致的，D 错误；

E．双星在周期性运动时，会使得到地球的距离发生周期性变化，故接收到的光频率会发生变化，E 正确。故选 BCE。

（2）(i) 如图，设距 c 点最近的振幅极大的点为 d 点，a 与 d 的距离为 $r_{1}$，b 与 d 的距离为 $r_{2}$，d 与 c 的距离为 $s$，

\onepicture[w=5.5cm]{\includesvg{figs/14_an}}

波长为 $\lambda$。则

$r_{2}-r_{1}=\lambda$

由几何关系有

$r_{1}=l-s$

$r_{2}^{2}=(r_{1}\sin60^\circ)^{2}+(l-r_{1}\cos60^\circ)^{2}$

联立各式并代入题给数据得

$\lambda=\frac{1}{4}l$

(ii) 波的频率为 $f$，设波的传播速度为 $v$，有

$v=f\lambda$

联立得 $v=\frac{fl}{4}$。
}

\end{enumerate}

\end{document}
